\documentclass[12pt]{amsart} 

\input{./latex-preamble/cm-preamble.tex}

%\graphicspath{ {./diagrams/} }

\DeclareMathOperator{\Div}{Div}
\DeclareMathOperator{\CaCl}{CaCl}
\DeclareMathOperator{\Pic}{Pic}

\newtheorem{lemma}{Lemma}

\begin{document}

{Hartshorne} \hfill {3-28-2025}

\bigskip\bigskip

\begin{center}

{\sc Chapter II.7, Week 9}

\end{center}

\begin{center}

  {Noah Parker} %\footnote{Collaborated with Kalev Martinson and Frank :).}}
    
\end{center}

\bigskip\bigskip

\begin{enumerate}
  \item[7.1.] {
      Let $(X,\O_X)$ be a locally ringed space, and let $f:\L\to \M$ be a surjective map of invertible sheaves on $X$.
      Show that $f$ is an isomorphism.
      [\emph{Hint:} Reduce to a question of modules over a local ring by looking at the stalks.]
    }
\end{enumerate}

\begin{enumerate}
  \item[7.2.] {
      Let $X$ be a scheme over a field $k$.
      Let $\L$ be an invertible sheaf on $X$, and let $\{s_0,\ldots,s_n\}$ and $\{t_0,\ldots,t_m\}$ be two sets of sections of $\L$, which generate the same subspace $V\subset \Gamma(X,\L)$, and which generate the sheaf $\L$ at every point.
      Suppose $n\leqs m$.
      Show that the corresponding morphisms $\phi:X\to \PP^n_k$ and $\psi:X\to \PP^m_k$ differ by a suitable linear projection $\PP^m-L\to \PP^n$, and an automorphism of $\PP^n$, where $L$ is a linear subspace of dimension $m-n-1$.
    }
\end{enumerate}

\begin{enumerate}
  \item[7.3.] {
      Let $\phi:\PP^n_k\to \PP^m_k$ be a morphism.
      Then:
      \begin{enumerate}
        \item either $\phi(\PP^n)=\{*\}$ or $m\geqs n$ and $\dim \phi(\PP^n)=n$;
        \item in the second case, $\phi$ can be obtained as the composition of (1) a $d$-uple embedding $\PP^n\to \PP^N$ for a uniquely determined $d\geqs 1$, (2) a linear projection $\PP^N-L\to \PP^m$, and (3) an automorphism of $\PP^m$.
      \end{enumerate}
    }
\end{enumerate}

\begin{enumerate}
  \item[7.5.] {
      Establish the following properties of ample and very ample invertible sheaves on a noetherian scheme $X$.
      $\L$, $\M$ will denote invertible sheaves, and for (d), (e) we assume furthermore that $X$ is of finite type over a noetherian ring $A$.
      \begin{enumerate}
        \item If $\L$ is ample and $\M$ is generated by global sections, then $\L\otimes \M$ is amples.
        \item If $\L$ is ample and $\M$ is arbitrary, then $\M\otimes \L^n$ is ample for $n\gg0$.
        \item If $\L$, $\M$ are both ample, so is $\L\otimes \M$.
        \item If $\L$ is very ample and $\M$ is generated by global sections, then $\L\otimes \M$ is very ample.
        \item If $\L$ is ample, then there is an $n_0>0$ such that $\L^n$ is very ample for all $n\geqs n_0$.
      \end{enumerate}
    }
\end{enumerate}

\begin{enumerate}
  \item[7.7.] {
      \emph{Some Rational Surfaces.}
      Let $X=\PP^2_k$, and let $|D|$ be the complete linear system of all divisors of degree 2 on $X$ (conics).
      $D$ corresponds to the invertible sheaf $\O(2)$, whose space of global sections has a basis $x^2,y^2,z^2,xy,xz,yz$, where $x,y,z$ are the homogneous coorrdinates of $X$.
      \begin{enumerate}
        \item The complete linear system $|D|$ gives an embedding of $\PP^2$ in $\PP^5$, whose image is the veronese surfaces (I, Ex. 2.13).
        \item Show that the subsystem defined by $x^2,y^2,z^2,y(x-z),(x-y)z$ gives a closed immersion of $X$ into $\PP^4$.
          The image is called the \emph{Veronese surface} in $\PP^4$. Cf. (IV, Ex. 3.11).
        \item Let $\fr{d}\subset |D|$ be the linear system of all conic passing through a fixed point $p$.
          Then $\fr d$ gives an imersion of $U=X-p$ into $\PP^4$.
          Furthermore, if we blow up $p$ to get a surface $\wtilde X$, then this map extends to give a closed immersion of $\wtilde X$ in $\PP^4$.
          Show that $\wtilde X$ is a surface of degree 3 in $\PP^4$, and that the lines in $X$ through $p$ are transformed into straight lines in $\wtilde X$ which do not meet.
          $\wtilde X$ is the union of all these lines, so we say $\wtilde X$ is a \emph{ruled surface} (V, 2.19.1).
      \end{enumerate}
    }
\end{enumerate}

\end{document}
